\section{Грассманова алгебра}%
\label{sec:Grassmannian_algebra}

Всюду в~этом подразделе векторы обозначаются полужирными буквами, а~греческими обозначаются числа (элементы поля или кольца).

Символом~$[m]$ обозначено множество
\[
	[m] := \set*{1, \dots, m}
	;
\]
символом $[m]^k$ обозначается $k$-ая декартова степень, то есть множество
\[
	[m]^k :=
	\underbrace{
		[m] \times \cdots \times [m]
	}_{k}
	:=
	\set{
		(i_1, \dots, i_k)
	}{
		\parbox{1.5in}{%
			\centering
			для всех~$\nu \in [k]$ %\par
			$i_\nu \in [m]$
		}
	}
	.
\]
Также символом~$\Delta^k$ обозначается множество наборов~$(i_1, \dots, i_k) \in [n]^k$, в~которых совпадают хотя бы два элемента~$i_\mu$ и~$i_\nu$.

Пусть $k \leq n$ и
\[
	I = (1 \leq i_1 < i_2 < \dots < i_k \leq n)
\]
--- набор из~$k$ упорядоченных по возрастанию элементов. Назовём \emph{дополнительным к~$I$ набором} набор
\[
	\hat{I} :=
	(1 \leq \hat{\imath}_1 < \hat{\imath}_2 < \dots < \hat{\imath}_{n-k} \leq n),
\]
состоящий из всех элементов из~$[n]$, не входящих в~$I$.
Для краткости будем обозначать~$\hat{I} = [n] \setminus I$.

Грассманово произведение обозначается~$\wedge$ \noatt{и~индуцируется тензорным произведением, факторизованным по соотношениям кососимметричности}. Иными словами, для любых векторов~$\vec{a}$, $\vec{b}$, $\vec{c}$ выполнены следующие равенства:
\begin{align*}
	\vec{a} \wedge \vec{a} &= 0,
	\\
	(\vec{a} + \vec{b}) \wedge \vec{c} &=
	\vec{a} \wedge \vec{c} +
	\vec{b} \wedge \vec{c},
	\\
	\vec{a} \wedge (\vec{b} + \vec{c}) &=
	\vec{a} \wedge \vec{b} +
	\vec{a} \wedge \vec{c},
	\\
	(\alpha \vec{a}) \wedge \vec{b} &=
	\alpha (\vec{a} \wedge \vec{b})
	,
	\\
	\vec{a} \wedge (\alpha \vec{b}) &=
	\alpha (\vec{a} \wedge \vec{b})
	.
\end{align*}
заметим, что из свойства~$\vec{a} \wedge \vec{a} = 0$ вытекает, что~$\vec{v} \wedge \vec{w} = -\vec{w} \wedge \vec{v}$ (достаточно подставить $\vec{a} := \vec{v} + \vec{w}$).

Одно из основных свойств грассмановых произведений вытекает из кососимметричности: если~$\sigma \in S_k$ --- перестановка, а~$(\vec{a}_1, \dots, \vec{a}_k)$ --- набор векторов из~$n$-мерноно пространства, то
\[
	\vec{a}_{\sigma(1)}
	\wedge
	\vec{a}_{\sigma(2)}
	\wedge
	\cdots
	\wedge
	\vec{a}_{\sigma(k)}
	=
	\sgn{\sigma} \cdot
	\vec{a}_{1}
	\wedge
	\vec{a}_{2}
	\wedge
	\cdots
	\wedge
	\vec{a}_{k}
	.
\]

Из этого следует удобное определение определителя.
\begin{definition}
	Пусть~$\ee = (\vec{e}_1, \dots \vec{e}_n)$ --- базис $n$-мерного пространства~$V$, а $\baa = (\vec{a}_1, \dots \vec{a}_n)$ --- набор из~$n$ векторов. (Обрати внимание, что векторов столько же, какова размерность пространства!)
	
	Если $\baa = \ee A$, то есть~$A$ --- матрица, состоящая из координат векторов~$\baa$ в~базисе~$\ee$, то~$\det{A}$ \emph{определяется} как множитель в~равенстве
	\[
		\vec{a}_1 \wcw \vec{a}_n
		=
		\det{A}
		\cdot
		\vec{e}_1 \wcw \vec{e}_n
		.
	\]
\end{definition}

\sloppy
Если хочется связать со знакомым определением, милости прошу. Пусть~$A = (a^i_j)$, где верхний индекс нумерует строки, а~нижний --- столбцы. Таким образом,
\[
	a_j =
	\Vect{a_j^1 \\ \dots \\ a_j^n}
\]
--- координаты вектора~$\vec{a}_j$, составляющие~$j$-ый столбец матрицы~$A$. Тогда
\begin{align*}
	\vec{a}_1 \wcw \vec{a}_n
	&=
	\left( 
		\sum_{i_1 = 1}^{n}
			a_1^{i_1}
			\vec{e_{i_1}}
	\right)
	\wcw
	\left( 
		\sum_{i_n = 1}^{n}
			a_n^{i_n}
			\vec{e_{i_n}}
	\right)
	\nlinea{=}
	\sum_{i_1 = 1}^{n}
	\cdots
	\sum_{i_n = 1}^{n}
		a_1^{i_1}
		\cdots
		a_n^{i_n}
		\cdot
		\vec{e}_{i_1}
		\wcw
		\vec{e}_{i_n}
	.
\end{align*}
Мы вынесли суммы и~числовые множители по линейности. Теперь заметим, что если в~наборе индексов~$(i_1, \dots, i_n)$ есть одинаковые, то
\[
	\vec{e}_{i_1} \wcw \vec{e}_{i_n} = 0
	.
\]
Поэтому останутся только слагаемые с~попарно разными индексами, и~сумму можно записать в~виде
\begin{multline*}
	\sum_{\sigma \in S_n}
	a_1^{\sigma(1)}
	\cdots
	a_n^{\sigma(n)}
	\cdot
	\vec{e}_{\sigma(1)}
	\wcw
	\vec{e}_{\sigma(n)}
	\nline{=}
	\sum_{\sigma \in S_n}
	a_1^{\sigma(1)}
	\cdots
	a_n^{\sigma(n)}
	\cdot
	\sgn{\sigma}
	\cdot
	\vec{e}_{1}
	\wcw
	\vec{e}_{n}
	\nline{=}
	\left(
		\sum_{\sigma \in S_n}
			\sgn{\sigma}
			\cdot
			a_1^{\sigma(1)}
			\cdots
			a_n^{\sigma(n)}
	\right)
	\cdot
	\vec{e}_{1}
	\wcw
	\vec{e}_{n}
	\nline{=}
	\det{A}
	\cdot
	\vec{e}_{1}
	\wcw
	\vec{e}_{n}
	.
\end{multline*}
Итак, определение определителя через грассманову алгебру совпадает со старым.

Спрашивается: а~зачем это всё? А~затем, чтобы можно было проще работать с~минорами, которые кодируют площади~$k$-мерных объёмов в~$n$-мерном пространстве. Если, например, 
\[
	\baa = (\vec{a}_1, \dots, \vec{a}_k) = \ee A
\]
--- набор из~$k < n$ векторов в~$n$-мерном пространстве, то определитель~$A$ взять не получится (матрица не квдаратная), зато получится взять грассманово произведение:
\begin{align*}
	\vec{a}_1 \wcw \vec{a}_k
	&=
	\left( 
		\sum_{i_1 = 1}^{n}
			a_1^{i_1}
			\vec{e}_{i_1}
	\right)
	\wcw
	\left( 
		\sum_{i_k = 1}^{n}
			a_k^{i_k}
			\vec{e}_{i_k}
	\right)
	\nlinea{=}
	\sum_{i_1 = 1}^{n}
	\cdots
	\sum_{i_k = 1}^{n}
		a_1^{i_1}
		\cdots
		a_k^{i_k}
	\cdot
		\vec{e}_{i_1}
		\wcw
		\vec{e}_{i_k}
	\nlinea{=}
	\sum_{(i_1, \dots, i_k) \in [n]^k \setminus \Delta}
		a_1^{i_1}
		\cdots
		a_k^{i_k}
	\cdot
		\vec{e}_{i_1}
		\wcw
		\vec{e}_{i_k}
	,
\end{align*}
где суммирование идёт по всем наборам~$(i_1, \dots, i_k) \in [n]$, попарно не совпадающими друг с~другом.

Частный случай этой ситуации, который для нас важен, --- это многомерный аналог векторного произведения, получающийся при~$k = n-1$. В~этом случае
\begin{align*}
	\vec{a}_1 \wcw \vec{a}_{n-1}
	&=
	\left( 
		\sum_{i_1 = 1}^{n}
			a_1^{i_1}
			\vec{e}_{i_1}
	\right)
	\wcw
	\left( 
		\sum_{i_{n-1} = 1}^{n}
			a_{n-1}^{i_{n-1}}
			\vec{e}_{i_{n-1}}
	\right)
	\nlinea{=}
	\sum_{i_1 = 1}^{n}
	\cdots
	\sum_{i_{n-1} = 1}^{n}
		a_1^{i_1}
		\cdots
		a_{n-1}^{i_{n-1}}
	\cdot
		\vec{e}_{i_1}
		\wcw
		\vec{e}_{i_{n-1}}
	\nlinea{=}
	\sum_{(i_1, \dots, i_{n-1}) \in [n] \setminus \Delta}
		a_1^{i_1}
		\cdots
		a_{n-1}^{i_{n-1}}
	\cdot
		\vec{e}_{i_1}
		\wcw
		\vec{e}_{i_{n-1}}
	.
\end{align*}
Суммирование ведётся по попарно не повторяющимся наборам чисел~$(i_1, \dots, i_{n-1})$. Но создать такой набор значит выбрать~$n-1$ число, входящее в~набор, из~$n$ чисел, или, что то же самое, \emph{выбрать одно число, которое не входит в~набор}. Поэтому суммирование можно вести по~\emph{одному} числу~$i$, \emph{не входящему в~грассманов моном}.
\sloppy

Каждое слагаемое
\[
	a_1^{i_1}
	\cdots
	a_k^{i_{n-1}}
	\cdot
	\vec{e}_{i_1}
	\wcw
	\vec{e}_{i_{n-1}}
\]
умножим на \emph{фиктивный} вектор~$\tta^i \te_i$, где $i$ --- единственное число, не входящее в~набор~$(i_1, \dots, i_{n-1})$. Тогда
\begin{multline*}
	\sum_{(i_1, \dots, i_{n-1}) \in [n] \setminus \Delta}
		a_1^{i_1}
		\cdots
		a_{n-1}^{i_{n-1}}
		\ta_i
	\cdot
		\te_i \wedge
		\vec{e}_{i_1}
		\wcw
		\vec{e}_{i_{n-1}}
	\nline{=}
	\sum_{(i_1, \dots, i_{n-1}) \in [n] \setminus \Delta}
		a_1^{i_1}
		\cdots
		a_{n-1}^{i_{n-1}}
		\ta_i
	\cdot
		\sgn(i_1, \cdots, i_{n-1})
	\cdot
		\te_i \wedge
		\vec{e}_1
		\wcw
		\widehat{\vec{e}}_i
		\wcw
		\vec{e}_n
	\nline{=}
	\sum_{(i_1, \dots, i_{n-1}) \in [n] \setminus \Delta}
		a_1^{i_1}
		\cdots
		a_{n-1}^{i_{n-1}}
		\ta_i
	\cdot
		\sgn(i_1, \cdots, i_{n-1})
		\cdot
		(-1)^{i-1}
	\cdot
		\vec{e}_1
		\wcw
		\te_i
		\wcw
		\vec{e}_n
	\nline{=}
	\sum_{i = 1}^n
		(-1)^{i-1}
		\ta_i
		\left( 
			\sum_{(i_1, \dots, i_{n-1}) \in S_{n-1}}
				a_1^{i_1}
				\cdots
				a_{n-1}^{i_{n-1}}
			\cdot
				\sgn(i_1, \cdots, i_{n-1})
		\right)
		\cdot
			\vec{e}_1
			\wcw
			\te_i
			\wcw
			\vec{e}_n
		.
\end{multline*}
Выражение в~скобках --- это определитель минора матрицы, в~которой вычеркнули первую строку и~$i$-ый столбец. А~сумма со знаками таких миноров, умноженных на~$\ta_i$, --- это определитель матрицы
\[
	\begin{pmatrix}
		\ta_1 & \ta_2 & \dots & \ta_n
		\\
		a_1^1 & a_1^2 & \dots & a_1^n
		\\
		\hdotsfor{4}
		\\
		a_{n-1}^1 & a_{n-1}^2 & \dots & a_{n-1}^n
	\end{pmatrix}
	.
\]
Обычно в~верхней строке такого определителя стоят формальные базисные векторы:
\[
	[\vec{a}_1, \dots, \vec{a}_{n-1}] :=
	\begin{determ}
		\vec{e}_1 & \vec{e}_2 & \dots & \vec{e}_n
		\\
		a_1^1 & a_1^2 & \dots & a_1^n
		\\
		\hdotsfor{4}
		\\
		a_{n-1}^1 & a_{n-1}^2 & \dots & a_{n-1}^n
	\end{determ}
	.
\]
Такой определитель естественно назвать \tdef{векторным произведением} векторов~$\vec{a}_1$, \dots, $\vec{a}_{n-1}$ в~$n$-мерном векторном пространстве~$V$.
\begin{remark}
	Определитель можно написать и~транспонированным:
	\[
		[\vec{a}_1, \dots, \vec{a}_{n-1}] =
		\begin{determ}
			\vec{e}_1 & a_1^1 & a_2^1 & \dots & a_{n-1}^1 \\
			\vec{e}_2 & a_1^2 & a_2^2 & \dots & a_{n-1}^2 \\
			\vec{e}_3 & a_1^3 & a_2^3 & \dots & a_{n-1}^3 \\
			\hdotsfor{5} \\
			\vec{e}_{n-1} & a_1^{n-1} & a_2^{n-1} & \dots & a_{n-1}^{n-1}
		\end{determ}
		.
	\]
	Это соответствует записи координат векторов в~столбик.
\end{remark}



% Обозначим~$[n] \setminus i$ набор чисел от~$1$ до~$n$, в~котором нет числа~$i$, а~для грассмановых мономов положим
% \[
% 	\vec{e}_{[n] \setminus i} :=
% 	\vec{e}_1
% 	\wcw
% 	\widehat{
% 		\vec{e}_i
% 	}
% 	\wcw
% 	\vec{e}_n
% \]
% (крышка означает, что множитель опущен).
% Аналогично обозначим
% \[
% 	a^{[n] \setminus i} :=
% 	a_1^{i_1} \cdots a_{n-1}^{i_{n-1}},
% \]
% где $(i_1, \dots, i_{n-1}) = [n] \setminus i$.
% Тогда
% \begin{align*}
% 	\vec{a}_1 \wcw \vec{a}_{n-1}
% 	&=
% 	\sum_{i = 1}^{n}
% 	\left( 
% 		\sum_{
% 			(i_1, \dots, i_{n-1})
% 			\in
% 			[n] \setminus i
% 		}
% 			a_1^{i_1}
% 			\cdots
% 			a_k^{i_{n-1}}
% 		\cdot
% 			\sgn(i_1, \dots, i_n)
% 			\vec{e}_{[n] \setminus i}
% 	\right)
% 	\nlinea{=}
% 	\sum_{i = 1}^{n}
% 		\sum_{
% 			(i_1, \dots, i_{n-1})
% 			\in
% 			[n] \setminus i
% 		}
% 	\left( 
% 		a_1^{i_1}
% 		\cdots
% 		a_k^{i_{n-1}}
% 		\cdot
% 		\sgn(i_1, \dots, i_n)
% 	\right)
% 	\vec{e}_{[n] \setminus i}
% 	.
% \end{align*}
% Если теперь приписать к~каждому моному
% $
% 	a_1^{i_1}
% 	\cdots
% 	a_k^{i_{n-1}}
% $
% \emph{фиктивный} векторный множитель~$\tvec{e}_i$, где $i \notin (i_1, \dots, i_{n-1})$, то
% \begin{align*}
% 	\vec{a}_1 \wcw \vec{a}_{n-1}
% 	&=
% 	\sum_{
% 		(i_1, \dots, i_{n-1})
% 		\in
% 		[n]^{n-1}\setminus\Delta
% 	}
% 		a_1^{i_1}
% 		\cdots
% 		a_k^{i_{n-1}}
% 		\cdot
% 		\sgn(i_1, \dots, i_n)




% \end{align*}

\subsection{Скалярные произведения}

В~грассмановой алгебре можно ввести скалярное произведение. Для мономов
$
	\vec{a}_1 \wcw \vec{a}_m
$
и
$
	\vec{b}_1 \wcw \vec{b}_m
$
по определению полагают
\[
	\spr{
		\vec{a}_1 \wcw \vec{a}_m
	}{
		\vec{b}_1 \wcw \vec{b}_m
	}
	:=
	\det
	\bigl(
		\spr{\vec{a}_i}{\vec{b}_j}
	\bigr)
	.
\]
Также вводят норму соотношением
\[
	\norm{
		\vec{a}_1 \wcw \vec{a}_m
	}^2
	:=
	\spr{
		\vec{a}_1 \wcw \vec{a}_m
	}{
		\vec{a}_1 \wcw \vec{a}_m
	}
	.
\]
Заметим, что
\[
	\spr{
		\vec{a}_1 \wcw \vec{a}_m
	}{
		\vec{a}_1 \wcw \vec{a}_m
	}
	=
	\det
	\bigl(
		\spr{\vec{a}_i}{\vec{a}_j}
	\bigr)
\]
--- определитель матрицы Грама набора векторов $\vec{a}_1$, \dots, $\vec{a}_m$.

Тогда $m$-мерная площадь параллелепипеда (в $n$-мерном пространстве), натянутого на $m$ векторов~$\vec{a}_1$, \dots, $\vec{a}_m$, равна
\[
	\vol_m(\vec{a}_1, \dots, \vec{a}_m) =
	\norm{
		\vec{a}_1 \wcw \vec{a}_m
	}.
\]
